\chapter{Sai Lầm Trong Suy Nghĩ}
\markboth{Sai Lầm Trong Suy Nghĩ}{Common Mistakes in Thinking}

\section{Nghĩ Mình Ngu}
\begin{truyen}{Nghĩ Mình Ngu}{Thinking You Are Stupid}
\chuhoa{C}{ó học sinh chỉ vì vài lần học chậm hoặc làm sai mà vội kết luận mình ngu.}
\textit{(Some students make a few mistakes or learn slowly and quickly conclude that they are stupid.)}

Kết luận ấy rất nặng, và thường không đúng.
\textit{(That conclusion is heavy, and often untrue.)}

Nó làm em mất động lực trước cả khi thật sự cố gắng đủ lâu.
\textit{(It kills motivation before you have truly tried for long enough.)}

Học chậm, hổng kiến thức, thiếu phương pháp hay đang mệt không có nghĩa là đầu óc em vô vọng.
\textit{(Learning slowly, having gaps, lacking method, or being tired does not mean your mind is hopeless.)}
\end{truyen}

\begin{baihoc}
Đừng biến một giai đoạn khó thành bản án cho cả con người mình.
\textit{(Do not turn a difficult phase into a life sentence for your whole self.)}
\end{baihoc}

\begin{ghinhoanh}
\textbf{Short English to remember:}

Slow progress is not proof of hopeless ability.

\end{ghinhoanh}

\begin{tuhocnhanh}
1. Vì sao câu “mình ngu” nguy hiểm? \textit{Why is the sentence “I am stupid” dangerous?}

2. Điều gì có thể làm em học kém tạm thời mà không phải vì ngu? \textit{What can hurt learning temporarily without proving stupidity?}

3. Em cần thay câu kết luận ấy bằng câu nào đúng hơn? \textit{What more accurate sentence should replace that conclusion?}
\end{tuhocnhanh}
\ngancach

\section{Nghĩ Mình Giỏi Quá}
\begin{truyen}{Nghĩ Mình Giỏi Quá}{Thinking You Are Already Too Good}
\chuhoa{T}{ự tin là tốt, nhưng tự tin quá mức dễ làm con người ngừng học.}
\textit{(Confidence is good, but overconfidence easily stops learning.)}

Khi nghĩ mình đã giỏi lắm rồi, em sẽ ít nghe góp ý, ít kiểm tra lại và ít chịu sửa lỗi.
\textit{(When you think you are already very good, you listen less, check less, and correct less.)}

Sai lầm này nguy hiểm vì nó làm em chậm tiến bộ trong khi vẫn tưởng mình ổn.
\textit{(This mistake is dangerous because it slows improvement while you still believe you are fine.)}

Người học tốt lâu dài thường vừa có tự tin, vừa có độ khiêm tốn đủ để tiếp tục học nữa.
\textit{(Strong lifelong learners usually hold both confidence and enough humility to keep learning.)}
\end{truyen}

\begin{baihoc}
Tự tin đúng giúp mình tiến lên; tự mãn làm mình đứng lại.
\textit{(Healthy confidence moves us forward; arrogance keeps us still.)}
\end{baihoc}

\begin{ghinhoanh}
\textbf{Short English to remember:}

Confidence grows best with humility.

\end{ghinhoanh}

\begin{tuhocnhanh}
1. Tự tin quá mức gây hại gì? \textit{What harm comes from overconfidence?}

2. Vì sao tự mãn làm việc học đứng lại? \textit{Why does arrogance stop learning?}

3. Em đang tự tin lành mạnh hay tự mãn? \textit{Are you healthily confident or becoming arrogant?}
\end{tuhocnhanh}
\ngancach

\section{So Sánh Với Người Khác}
\begin{truyen}{So Sánh Với Người Khác}{Comparing Yourself to Others}
\chuhoa{S}{o sánh là phản xạ rất tự nhiên, nhất là trong môi trường học đường.}
\textit{(Comparison is a natural reflex, especially in school.)}

Nhưng nếu sống bằng việc nhìn sang người khác quá nhiều, em sẽ khó nhìn rõ đường của mình.
\textit{(But if you live by watching others too much, you will struggle to see your own path.)}

So sánh có thể khiến em tự ti hoặc kiêu ngạo, mà cả hai đều không giúp học tốt hơn.
\textit{(Comparison can produce inferiority or arrogance, and neither helps learning.)}

Điều hữu ích hơn là nhìn xem hôm nay mình khá hơn chính mình hôm trước chưa.
\textit{(It is more useful to ask whether you are better today than before.)}
\end{truyen}

\begin{baihoc}
So sánh quá nhiều làm mình mệt, nhưng so mình với phiên bản cũ của mình thường làm mình lớn lên.
\textit{(Too much comparison exhausts us, but comparing ourselves with our earlier self often helps us grow.)}
\end{baihoc}

\begin{ghinhoanh}
\textbf{Short English to remember:}

Compare less with others, more with your past self.

\end{ghinhoanh}

\begin{tuhocnhanh}
1. So sánh quá nhiều gây ra điều gì? \textit{What does too much comparison create?}

2. Vì sao nên nhìn lại chính mình nhiều hơn? \textit{Why should you focus more on your own growth?}

3. Em đã tiến bộ hơn chính mình cũ ở điểm nào? \textit{In what way have you improved from your past self?}
\end{tuhocnhanh}
\ngancach

\section{Sợ Bị Chê}
\begin{truyen}{Sợ Bị Chê}{Fear of Criticism}
\chuhoa{N}{ỗi sợ bị chê làm nhiều học sinh thu mình lại trước cả khi thử.}
\textit{(Fear of criticism causes many students to shrink back before even trying.)}

Các em sợ hỏi, sợ phát biểu, sợ nộp bài, sợ lộ chỗ còn yếu.
\textit{(They fear asking, speaking up, submitting work, or exposing weakness.)}

Nhưng càng sợ chê và càng tránh, em càng ít cơ hội để khá hơn.
\textit{(But the more you avoid criticism, the fewer chances you have to improve.)}

Không phải lời chê nào cũng đúng, nhưng nỗi sợ quá lớn luôn làm con người nhỏ đi.
\textit{(Not every criticism is valid, but overwhelming fear always shrinks a person.)}
\end{truyen}

\begin{baihoc}
Đừng để nỗi sợ lời người khác lớn hơn nhu cầu học và lớn lên của chính mình.
\textit{(Do not let fear of others' words outweigh your need to learn and grow.)}
\end{baihoc}

\begin{ghinhoanh}
\textbf{Short English to remember:}

Fear of criticism should not rule growth.

\end{ghinhoanh}

\begin{tuhocnhanh}
1. Sợ bị chê khiến học sinh bỏ lỡ điều gì? \textit{What does fear of criticism make students lose?}

2. Vì sao tránh mãi không giúp mình mạnh hơn? \textit{Why does constant avoidance not make us stronger?}

3. Em đang né điều gì vì sợ người khác nói? \textit{What are you avoiding because of others' opinions?}
\end{tuhocnhanh}
\ngancach

\section{Nghĩ Thất Bại Là Hết}
\begin{truyen}{Nghĩ Thất Bại Là Hết}{Thinking Failure Means the End}
\chuhoa{M}{ột lần thất bại có thể làm em buồn, nhưng nó không phải dấu chấm hết.}
\textit{(One failure can hurt, but it is not the final end.)}

Sai lầm ở đây là biến một lần vấp thành kết luận rằng không còn cơ hội nữa.
\textit{(The mistake is turning one setback into the conclusion that there is no hope left.)}

Nhiều người tiến bộ tốt chính nhờ họ dám nhìn lại thất bại như dữ liệu để sửa.
\textit{(Many people improve precisely because they treat failure as data for correction.)}

Thất bại đau, nhưng nếu nhìn đúng, nó còn dạy được rất nhiều.
\textit{(Failure hurts, but if viewed rightly, it also teaches a great deal.)}
\end{truyen}

\begin{baihoc}
Thất bại là một đoạn đường khó, không phải bản đồ của cả đời em.
\textit{(Failure is a hard stretch of road, not the map of your whole life.)}
\end{baihoc}

\begin{ghinhoanh}
\textbf{Short English to remember:}

Failure is a lesson, not a final identity.

\end{ghinhoanh}

\begin{tuhocnhanh}
1. Vì sao không nên nghĩ thất bại là hết? \textit{Why should failure not be seen as the end?}

2. Thất bại có thể dạy điều gì? \textit{What can failure teach?}

3. Em cần nhìn lại thất bại nào bằng con mắt khác? \textit{Which failure do you need to re-examine differently?}
\end{tuhocnhanh}
\ngancach

\section{Nghĩ Học Là Khổ}
\begin{truyen}{Nghĩ Học Là Khổ}{Thinking Learning Is Only Suffering}
\chuhoa{K}{hi nghĩ học chỉ là khổ, học sinh sẽ bước vào việc học với tâm thế chống đối.}
\textit{(When students think learning is only suffering, they approach it with resistance.)}

Đúng là học có lúc mệt, khó và áp lực.
\textit{(It is true that learning can be tiring, hard, and stressful.)}

Nhưng học cũng là cách mở đầu óc, mở cơ hội và mở cả cách hiểu đời.
\textit{(But learning also opens the mind, opportunities, and ways of understanding life.)}

Nếu chỉ nhìn mặt khổ, em sẽ bỏ lỡ mặt đáng giá của nó.
\textit{(If you see only the suffering, you miss what makes it valuable.)}
\end{truyen}

\begin{baihoc}
Học không phải lúc nào cũng dễ chịu, nhưng gọi nó chỉ là khổ thì quá hẹp.
\textit{(Learning is not always comfortable, but calling it nothing but suffering is too narrow.)}
\end{baihoc}

\begin{ghinhoanh}
\textbf{Short English to remember:}

Learning is difficult, but it is also meaningful.

\end{ghinhoanh}

\begin{tuhocnhanh}
1. Vì sao nghĩ học là khổ làm mình học kém hơn? \textit{Why does seeing learning only as suffering weaken learning?}

2. Ngoài áp lực, việc học còn cho mình điều gì? \textit{What does learning give beyond pressure?}

3. Em có đang nhìn việc học quá nặng một phía không? \textit{Are you seeing learning too one-sidedly?}
\end{tuhocnhanh}
\ngancach

\section{Nghĩ Mình Không Thay Đổi Được}
\begin{truyen}{Nghĩ Mình Không Thay Đổi Được}{Thinking You Cannot Change}
\chuhoa{C}{ó em nói: “Tính em thế rồi,” như thể bản thân là thứ đã đóng khung.}
\textit{(Some students say, “That is just how I am,” as though the self were permanently fixed.)}

Nhưng con người thay đổi rất nhiều qua cách nghĩ, cách tập và cách sống mỗi ngày.
\textit{(Yet people change greatly through how they think, practice, and live each day.)}

Không phải điều gì cũng đổi ngay, nhưng rất nhiều thứ có thể tốt dần lên.
\textit{(Not everything changes immediately, but many things can improve over time.)}

Nghĩ mình không thay đổi được là tự đóng cửa trước khi thử mở nó.
\textit{(Thinking you cannot change is closing the door before even trying to open it.)}
\end{truyen}

\begin{baihoc}
Thay đổi thường chậm, nhưng điều đó khác hẳn với chuyện không thể thay đổi.
\textit{(Change is often slow, but that is very different from being impossible.)}
\end{baihoc}

\begin{ghinhoanh}
\textbf{Short English to remember:}

Slow change is still real change.

\end{ghinhoanh}

\begin{tuhocnhanh}
1. Vì sao câu “em không đổi được đâu” là nguy hiểm? \textit{Why is the sentence “I cannot change” dangerous?}

2. Con người thay đổi bằng những gì? \textit{Through what does a person change?}

3. Điều nào của em có thể tốt lên nếu kiên trì tập? \textit{What part of you could improve with steady practice?}
\end{tuhocnhanh}
\ngancach

\section{Nghĩ Người Khác Luôn Giỏi Hơn}
\begin{truyen}{Nghĩ Người Khác Luôn Giỏi Hơn}{Thinking Others Are Always Better}
\chuhoa{N}{hiều học sinh nhìn thấy điểm mạnh của người khác nhưng lại nhìn quá kỹ điểm yếu của mình.}
\textit{(Many students notice others' strengths while staring too hard at their own weaknesses.)}

Từ đó, các em nghĩ người khác lúc nào cũng hơn mình hết.
\textit{(From there, they assume everyone else is better in every way.)}

Thật ra ai cũng có chỗ mạnh, chỗ yếu và những nỗi lo không nói ra.
\textit{(In reality, everyone has strengths, weaknesses, and silent worries.)}

Nhìn người khác quá cao còn tự nhìn mình quá thấp chỉ làm em khó phát huy phần tốt của mình.
\textit{(Overestimating others and underestimating yourself only prevents your own strengths from growing.)}
\end{truyen}

\begin{baihoc}
Em không cần phủ nhận điểm mạnh của người khác, nhưng cũng đừng xóa sạch giá trị của mình.
\textit{(You do not need to deny others' strengths, but do not erase your own worth either.)}
\end{baihoc}

\begin{ghinhoanh}
\textbf{Short English to remember:}

Others may be strong, but that does not erase your value.

\end{ghinhoanh}

\begin{tuhocnhanh}
1. Vì sao học sinh dễ nghĩ người khác luôn giỏi hơn? \textit{Why do students easily think others are always better?}

2. Cách nhìn này làm mình thiệt ở đâu? \textit{How does this way of thinking hurt you?}

3. Em có điểm mạnh nào mình hay quên mất? \textit{What strength of yours do you often forget?}
\end{tuhocnhanh}
\ngancach

\section{Nghĩ Mình Không Có Tương Lai}
\begin{truyen}{Nghĩ Mình Không Có Tương Lai}{Thinking You Have No Future}
\chuhoa{K}{hi hiện tại rối, học hành chật vật, nhiều em rất dễ nghĩ phía trước mình chẳng có gì sáng cả.}
\textit{(When the present feels messy and learning is difficult, many students easily think nothing bright lies ahead.)}

Nhưng tương lai không được quyết định trọn vẹn chỉ bởi một giai đoạn khó.
\textit{(But the future is not fully determined by one difficult phase.)}

Nó còn thay đổi theo cách em tiếp tục hay bỏ cuộc, học thêm hay mặc kệ, sửa mình hay tự kết thúc sớm.
\textit{(It also shifts according to whether you continue or quit, learn more or give up, improve yourself or close your story too early.)}

Tương lai có thể mờ, nhưng mờ không có nghĩa là không có.
\textit{(The future may be unclear, but unclear does not mean absent.)}
\end{truyen}

\begin{baihoc}
Đừng dùng một quãng tối để kết luận rằng trước mặt mình không còn đường nào sáng nữa.
\textit{(Do not use one dark stretch to conclude that no bright road remains ahead.)}
\end{baihoc}

\begin{ghinhoanh}
\textbf{Short English to remember:}

An unclear future is not the same as no future.

\end{ghinhoanh}

\begin{tuhocnhanh}
1. Vì sao nhiều em nghĩ mình không có tương lai? \textit{Why do many students think they have no future?}

2. Tương lai thay đổi theo những gì? \textit{What can still change the future?}

3. Khi thấy mù mờ, em cần bám vào điều gì trước mắt? \textit{When things feel unclear, what should you hold onto first?}
\end{tuhocnhanh}
\ngancach

\section{Nghĩ Cố Gắng Là Vô Ích}
\begin{truyen}{Nghĩ Cố Gắng Là Vô Ích}{Thinking Effort Is Useless}
\chuhoa{Đ}{ây là một trong những ý nghĩ làm người học gục nhanh nhất.}
\textit{(This is one of the thoughts that defeats learners the fastest.)}

Khi tin cố gắng là vô ích, em sẽ không còn lý do để bắt đầu hay kiên trì.
\textit{(When you believe effort is useless, you lose the reason to start or persist.)}

Dĩ nhiên không phải cứ cố là được ngay, nhưng thiếu cố gắng thì gần như chắc chắn không có gì khá lên.
\textit{(Of course effort does not guarantee instant success, but without effort improvement is almost impossible.)}

Cố gắng không luôn cho kết quả nhanh, nhưng nó vẫn là cây cầu thật nhất nối hiện tại với khả năng tốt hơn.
\textit{(Effort does not always produce quick results, but it remains the truest bridge between the present and a better possibility.)}
\end{truyen}

\begin{baihoc}
Cố gắng không phải phép màu, nhưng từ bỏ gần như luôn là cách chắc nhất để không đổi được gì.
\textit{(Effort is not magic, but giving up is almost always the surest way to change nothing.)}
\end{baihoc}

\begin{ghinhoanh}
\textbf{Short English to remember:}

Effort may be slow, but quitting is emptier.

\end{ghinhoanh}

\begin{tuhocnhanh}
1. Vì sao nghĩ cố gắng là vô ích rất nguy hiểm? \textit{Why is thinking effort is useless so dangerous?}

2. Cố gắng không bảo đảm điều gì, nhưng vẫn cần vì sao? \textit{What does effort not guarantee, and why is it still necessary?}

3. Em đang muốn bỏ điều gì vì thấy cố cũng vô ích? \textit{What are you tempted to quit because effort feels pointless?}
\end{tuhocnhanh}
\ngancach